% !TEX root = ../enac_project_fr.tex
%%%%%%%%%%%%%%%%%%%%%%%%%%%%%%%%%%%%%%%%%%%%%%%%%%%%%%%%%%%%%%%%%%%%
% Chapter 4 - Implémentation dans Orekit
%%%%%%%%%%%%%%%%%%%%%%%%%%%%%%%%%%%%%%%%%%%%%%%%%%%%%%%%%%%%%%%%%%%%

\chapter{Implémentation dans Orekit}
\label{implementation_orekit}

\section{Architecture logicielle}

L'ensemble des algorithmes développés au cours de ce stage a été implémenté
en Java en s'appuyant sur la bibliothèque \textit{Orekit}. Celle-ci fournit
des mécanismes génériques pour la représentation des états, la propagation,
la modélisation des mesures et leur exploitation dans des processus
d'estimation.

L'objectif de l'implémentation n'est pas de reproduire indépendamment
l'ensemble de ces mécanismes, mais de les adapter au problème spécifique du
positionnement GNSS d'un récepteur au sol. L'architecture développée repose
ainsi sur une séparation entre l'acquisition et la préparation des données,
la construction des mesures, leur modélisation et leur résolution.

La Figure~\ref{fig:architecture_logicielle} présente l'organisation
générale de l'application.

\begin{figure}[htbp]
    \centering
    \rule{0.9\linewidth}{6cm} % TODO : diagramme UML simplifié
    \caption{Architecture générale du logiciel développé.}
    \label{fig:architecture_logicielle}
\end{figure}

La chaîne de traitement peut être résumée par les étapes suivantes :

\begin{enumerate}
    \item lecture et préparation des données RINEX ;
    \item sélection et regroupement des observations par époque et par
    satellite ;
    \item construction des objets de mesure Orekit ;
    \item ajout des modèles de correction nécessaires ;
    \item définition des paramètres susceptibles d'être estimés ;
    \item résolution du problème par moindres carrés pondérés ou par filtre
    de Kalman.
\end{enumerate}

Deux méthodes d'estimation ont été développées :

\begin{itemize}
    \item un solveur par moindres carrés pondérés
    (\textit{Weighted Least Squares}, WLS), principalement utilisé pour le
    calcul de solutions SPP ;
    \item un solveur fondé sur un filtre de Kalman étendu (EKF), utilisé pour
    l'estimation récursive des paramètres et constituant la base de
    l'implémentation du PPP.
\end{itemize}

Cette organisation permet de modifier indépendamment les modèles physiques
et la méthode d'estimation. Il est ainsi possible, par exemple, de modifier
le modèle troposphérique ou le nombre de paramètres estimés sans avoir à
reconstruire manuellement l'ensemble du système de résolution.

\section{Construction des observations GNSS}

Les données extraites des fichiers RINEX doivent être transformées en une
représentation adaptée aux mécanismes d'estimation d'Orekit. Cette étape
constitue l'interface entre les données GNSS et le modèle de mesure utilisé
par les solveurs.

Le traitement est réalisé indépendamment pour chaque époque d'observation.
Les observations disponibles dans le fichier RINEX sont d'abord parcourues
afin d'identifier les satellites observés et les mesures associées.

Une première sélection permet d'éliminer les observations ne satisfaisant
pas les critères définis pour l'expérience. Les principaux critères
appliqués sont :

\begin{itemize}
    \item la constellation GNSS considérée ;
    \item la disponibilité des observations nécessaires au traitement ;
    \item l'élévation minimale du satellite ;
    \item le rapport signal sur bruit minimal ;
    \item l'état de santé du satellite.
\end{itemize}

Pour chaque satellite retenu, les éphémérides correspondantes sont ensuite
recherchées dans les données de navigation. Elles permettent notamment de
déterminer l'état du satellite à l'instant d'émission du signal.

Les informations nécessaires au traitement sont alors regroupées dans une
représentation intermédiaire comprenant notamment :

\begin{itemize}
    \item l'identifiant du satellite ;
    \item sa constellation ;
    \item les observations de code et de phase ;
    \item sa position et sa vitesse ;
    \item son biais d'horloge ;
    \item les informations nécessaires aux modèles de correction.
\end{itemize}

Cette étape permet de dissocier la lecture des fichiers RINEX de la
construction proprement dite des mesures. Les données nécessaires à une
observation sont ainsi rassemblées avant leur transmission aux mécanismes
d'estimation d'Orekit.

\subsection{Création des objets de mesure}

Les observations préparées précédemment sont finalement représentées par
des objets de mesure Orekit. Deux classes principales sont utilisées :

\begin{itemize}
    \item \texttt{OneWayGNSSRange}, pour les observations de pseudo-distance ;
    \item \texttt{OneWayGNSSPhase}, pour les observations de phase de la
    porteuse.
\end{itemize}

Ces classes implémentent l'interface générique
\texttt{ObservedMeasurement}, qui fournit notamment les mécanismes
nécessaires pour associer une valeur observée à une valeur théorique calculée
par un modèle de mesure. :contentReference[oaicite:1]{index=1}

Les objets ainsi créés constituent donc le support commun utilisé par les
deux méthodes de résolution développées dans ce travail.

\section{Implémentation des modèles de correction}

Une fois les objets de mesure créés, ceux-ci doivent être complétés par les
différents modèles physiques intervenant dans les équations d'observation.
Cette étape constitue une partie importante du développement réalisé au
cours du stage.

Orekit permet d'associer à une mesure différents composants chargés de
modifier sa valeur théorique. Cette architecture permet de traiter
séparément les différentes contributions physiques au modèle de mesure.

Certaines corrections nécessaires au positionnement GNSS étaient déjà
disponibles dans Orekit. Elles ont donc pu être directement réutilisées dans
l'implémentation. D'autres, en revanche, n'étaient pas disponibles sous une
forme adaptée au positionnement d'un récepteur au sol et ont nécessité le
développement de nouveaux composants.

\section{Adaptation et extension des modèles Orekit}
\label{sec:adaptation_orekit}

L'un des objectifs principaux du développement a consisté à exploiter au maximum
les mécanismes de modélisation et d'estimation déjà disponibles dans Orekit,
tout en les complétant lorsque ceux-ci ne répondaient pas directement aux
besoins du positionnement GNSS considéré.

Cette démarche repose sur une distinction entre deux catégories de
fonctionnalités. Certaines corrections et certains mécanismes nécessaires au
positionnement sont déjà implémentés dans Orekit et peuvent être directement
réutilisés. D'autres nécessitent au contraire une adaptation ou un
développement spécifique, notamment lorsqu'ils doivent être appliqués aux
objets de mesure \texttt{OneWayGNSSRange} et \texttt{OneWayGNSSPhase} utilisés
dans cette étude.

Le tableau~\ref{tab:orekit_corrections} synthétise les principales corrections
considérées et indique leur disponibilité initiale dans Orekit ainsi que, le
cas échéant, le composant développé dans le cadre du stage.

\begin{table}[H]
\centering
\caption{Disponibilité et extensions développées pour les modèles de correction.}
\label{tab:orekit_corrections}
\small
\begin{tabular}{p{5.0cm} c p{6.0cm}}
\hline
\textbf{Correction / modèle} &
\textbf{Disponible dans Orekit} &
\textbf{Extension développée} \\
\hline

Horloge satellite &
Oui &
\texttt{AbstractSatelliteClockOneWayGNSSModifier},
\texttt{SatelliteClockOneWayGNSSRangeModifier},
\texttt{SatelliteClockOneWayGNSSPhaseModifier}
\\

Correction d'antenne L1 -- phase &
Partiellement &
\texttt{L1AntennaPhaseModifier}
\\

Correction d'antenne L1 -- pseudo-distance &
Partiellement &
\texttt{L1AntennaRangeModifier}
\\

Correction d'antenne bi-fréquence -- phase &
Non directement adaptée &
\texttt{IonosphereFreeAntennaPhaseModifier}
\\

Correction d'antenne bi-fréquence -- pseudo-distance &
Non directement adaptée &
\texttt{IonosphereFreeAntennaRangeModifier}
\\

Correction ionosphérique -- phase &
Non pour le modèle utilisé &
\texttt{IonosphericDelayOneWayGNSSPhaseModifier}
\\

Correction ionosphérique -- pseudo-distance &
Non pour le modèle utilisé &
\texttt{IonosphericDelayOneWayGNSSRangeModifier}
\\

Correction troposphérique -- phase &
Non pour l'application considérée &
\texttt{TroposphericDelayOneWayGNSSPhaseModifier}
\\

Correction troposphérique -- pseudo-distance &
Non pour l'application considérée &
\texttt{TroposphericDelayOneWayGNSSRangeModifier}
\\

Wind-up de phase &
Non directement adapté &
\texttt{WindUpOneWayGNSSPhaseModifier}
\\

Biais inter-systèmes Galileo &
Partiellement &
\texttt{GalileoIsbModifier}
\\

Détection des sauts de cycle &
Non &
\texttt{RtkSlipDetector}
\\

Filtrage des innovations &
Non &
\texttt{RtkMaxInnoFilter}
\\

\hline
\end{tabular}
\end{table}

\subsection{Corrections existantes réutilisées}

Les mécanismes existants dans Orekit permettent notamment de prendre en
compte certaines corrections géométriques, relativistes et liées aux
antennes. Leur intégration sous forme de composants indépendants permet de
conserver la structure modulaire du modèle de mesure.

L'utilisation de ces mécanismes permet également d'éviter de reproduire dans
le code applicatif des modèles physiques déjà implémentés et validés dans
Orekit.

\subsection{Développement des modèles atmosphériques}

Les modèles atmosphériques constituent l'une des principales extensions
réalisées au cours du stage.

Les modèles nécessaires à la modélisation de la troposphère et de
l'ionosphère n'étaient pas directement disponibles sous la forme requise
pour le traitement développé. Il a donc été nécessaire de concevoir des
composants permettant d'intégrer ces corrections directement aux mesures
\texttt{OneWayGNSSRange} et \texttt{OneWayGNSSPhase}.

Les modèles développés reprennent les formulations théoriques présentées au
chapitre précédent. Leur implémentation permet notamment de considérer
différentes configurations du modèle troposphérique, avec ou sans estimation
de paramètres supplémentaires.

Cette organisation permet ainsi de conserver une séparation claire entre le
modèle physique de propagation et le mécanisme de résolution utilisé pour
estimer les paramètres.

\subsection{Paramètres estimables}

Certaines grandeurs intervenant dans les équations d'observation ne sont pas
considérées comme des corrections déterministes. Elles sont au contraire
introduites comme des paramètres susceptibles d'être estimés avec les
observations.

Orekit fournit pour cela la classe \texttt{ParameterDriver}. Un
\texttt{ParameterDriver} permet notamment de définir la valeur d'un
paramètre, son échelle de normalisation, son caractère estimable et les
informations nécessaires à son intégration dans le processus d'estimation.
La documentation d'Orekit indique notamment que les paramètres estimés par
le filtre de Kalman sont pilotés par ces objets, et distingue les paramètres
orbitaux, de propagation et de mesure. :contentReference[oaicite:2]{index=2}

Dans l'application développée, ce mécanisme est notamment utilisé pour
représenter :

\begin{itemize}
    \item le biais d'horloge du récepteur ;
    \item les ambiguïtés de phase ;
    \item les paramètres troposphériques lorsque ceux-ci sont estimés ;
    \item les autres paramètres de mesure sélectionnés selon la configuration
    utilisée.
\end{itemize}

Cette architecture permet de modifier la dimension du problème d'estimation
sans modifier explicitement les matrices utilisées par le solveur.

\section{Résolution}
\label{sec:resolution_orekit}

Une fois les observations construites, les corrections appliquées et les
paramètres estimables définis, l'ensemble peut être transmis au mécanisme
de résolution.

L'un des choix importants réalisés au cours du développement a consisté à
réutiliser les mécanismes génériques d'estimation d'Orekit plutôt qu'à
réimplémenter directement les algorithmes de moindres carrés et de Kalman.

Cette approche présente deux avantages principaux. D'une part, elle permet
de bénéficier des implémentations numériques existantes dans Orekit et dans
la bibliothèque mathématique Hipparchus. D'autre part, elle évite de coder
manuellement une structure matricielle différente pour chaque configuration
du problème.

En effet, la dimension du vecteur d'état et la structure des dérivées
peuvent évoluer selon les paramètres sélectionnés pour l'estimation. Par
exemple, l'ajout d'un paramètre troposphérique ou d'une ambiguïté de phase
modifie directement le nombre d'inconnues du problème.

Les objets de mesure et les \texttt{ParameterDriver} permettent de transmettre
cette information au moteur d'estimation, qui peut alors construire
dynamiquement les grandeurs nécessaires à la résolution.

\begin{table}[htbp]
\centering
\caption{Principales classes Orekit utilisées dans l'implémentation.}
\label{tab:orekit_classes_estimation}
\begin{tabular}{ll}
\hline
\textbf{Classe / objet} & \textbf{Rôle dans l'implémentation} \\
\hline
\texttt{OneWayGNSSRange} & Représentation d'une mesure de pseudo-distance \\
\texttt{OneWayGNSSPhase} & Représentation d'une mesure de phase \\
\texttt{ParameterDriver} & Gestion des paramètres estimables \\
\texttt{AmbiguityDriver} & Gestion des ambiguïtés de phase \\
\texttt{EstimatedMeasurement} & Évaluation théorique d'une mesure \\
\texttt{BatchLSEstimator} & Résolution par moindres carrés \\
\texttt{KalmanEstimator} & Résolution séquentielle par filtre de Kalman \\
\hline
\end{tabular}
\end{table}

\subsection{Résolution par moindres carrés pondérés}

La résolution SPP repose sur une estimation par moindres carrés pondérés.
Dans l'implémentation, les observations sont préalablement associées à leur
modèle de mesure et à leurs incertitudes.

Une fois cette configuration réalisée, la résolution est confiée au
\texttt{BatchLSEstimator} d'Orekit. Celui-ci s'appuie sur les outils de
résolution de problèmes de moindres carrés de la bibliothèque Hipparchus.
L'architecture d'Orekit prévoit notamment un modèle de moindres carrés
construit à partir des \texttt{ObservedMeasurement} et des paramètres
estimés. :contentReference[oaicite:3]{index=3}

Dans cette implémentation, l'appel au solveur est donc relativement simple :
la complexité est principalement reportée en amont, lors de la construction
des observations et de leur modèle.

Cette organisation évite notamment de construire explicitement les matrices
du système pour chaque configuration. Si le nombre de paramètres estimés
évolue, la structure correspondante est automatiquement prise en compte par
le mécanisme d'estimation.

L'optimisation numérique utilisée repose sur
\texttt{LevenbergMarquardtOptimizer} de Hipparchus :

\begin{center}
\texttt{org.hipparchus.optim.nonlinear.vector.leastsquares.}\\
\texttt{LevenbergMarquardtOptimizer}.
\end{center}

Le développement porte ainsi principalement sur la définition correcte des
mesures, des modèles physiques et des paramètres, plutôt que sur la
réimplémentation de l'algorithme numérique de résolution.

\subsection{Résolution par filtre de Kalman}

Le second mode de résolution repose sur le filtre de Kalman étendu fourni
par Orekit. Dans la bibliothèque, le \texttt{KalmanEstimator} est initialement
destiné aux problèmes de détermination orbitale et utilise un
\texttt{PropagatorBuilder} pour définir la trajectoire de référence. Les
paramètres estimés sont pilotés par des \texttt{ParameterDriver}. :contentReference[oaicite:4]{index=4}

L'application de ce mécanisme au positionnement d'un récepteur fixe nécessite
cependant d'adapter la représentation de l'état et son évolution temporelle.
En particulier, le mécanisme de filtrage repose sur un propagateur permettant
de fournir l'état à l'époque suivante.

Or, une station GNSS fixe ne suit pas une trajectoire orbitale et ne nécessite
pas de propagation dynamique au sens classique de la détermination orbitale.
Un propagateur spécifique, nommé \texttt{StaticReceiverPropagator}, a donc
été développé afin de représenter l'état d'un récepteur stationnaire tout en
fournissant à l'estimateur les informations nécessaires à son fonctionnement.

Cette adaptation permet de réutiliser le mécanisme de Kalman d'Orekit pour
un problème dont la formulation physique diffère de son application
initiale. Le principe consiste à conserver l'architecture générique du moteur
d'estimation tout en remplaçant la dynamique orbitale par une évolution
adaptée à une station fixe.

Le filtre exploite ensuite les observations construites précédemment ainsi
que les paramètres sélectionnés pour l'estimation. Les incertitudes des
mesures de code et de phase définissent leur pondération, tandis que les
bruits de processus associés aux paramètres estimés permettent de définir
leur évolution temporelle.

La structure du filtre est ainsi déterminée par les paramètres effectivement
sélectionnés. Une estimation comprenant uniquement la position et l'horloge
ne possède pas la même dimension qu'une estimation incluant également les
ambiguïtés de phase ou des paramètres troposphériques.

Comme pour les moindres carrés, aucune construction manuelle des matrices
du filtre n'est donc nécessaire dans le code applicatif. Les paramètres,
leurs échelles et leurs incertitudes sont fournis au mécanisme d'estimation,
qui assure ensuite la construction et la mise à jour des grandeurs
nécessaires au filtrage. La documentation d'Orekit précise notamment que les
états, covariances et matrices de mesure utilisés par le filtre sont
manipulés par le moteur de Kalman et par Hipparchus, avec une normalisation
associée aux paramètres estimés et aux incertitudes de mesure. :contentReference[oaicite:5]{index=5}

\section{Synthèse de la chaîne d'implémentation}

L'architecture finalement développée peut être résumée par la chaîne
suivante :

\begin{center}
\textit{RINEX}
$\rightarrow$
\textit{sélection des observations}
$\rightarrow$
\textit{objets de mesure Orekit}
$\rightarrow$
\textit{modèles de correction}
$\rightarrow$
\textit{paramètres estimables}
$\rightarrow$
\textit{WLS ou EKF}
$\rightarrow$
\textit{solution GNSS}.
\end{center}

Cette architecture permet de distinguer clairement les trois niveaux du
développement. Les données GNSS sont d'abord transformées en objets de mesure,
les modèles physiques sont ensuite associés à ces mesures, puis le problème
ainsi défini est transmis à un moteur d'estimation générique.

L'essentiel du travail de développement a ainsi porté sur l'adaptation des
mécanismes génériques d'Orekit au contexte spécifique du positionnement GNSS
au sol, notamment par l'implémentation des modèles atmosphériques et du
propagateur associé à une station fixe. Les mécanismes numériques de
résolution sont quant à eux réutilisés à travers les outils d'estimation
d'Orekit et de Hipparchus.